% Recto verso : pages impaires à droite, pages paires à gauche.
\documentclass[a4paper,12pt,twoside]{article}

% --------------------------------------------------
% Language and typography — configuration for pdfLaTeX
% --------------------------------------------------

% Encodage des fontes utilisées pour composer le texte.
%
% T1, aussi appelé encodage Cork, définit 256 positions de caractères
% dans une fonte, notamment des lettres accentuées européennes.
% Ce n'est ni une police particulière, ni l'encodage du fichier source.
%
% Avec l'ancien encodage OT1, certains caractères accentués sont
% construits en superposant un accent et une lettre.
% T1 permet de les utiliser comme des caractères à part entière,
% ce qui facilite notamment la césure des mots accentués.
%
% Exemple : écrire "café" ou le nom d'un auteur avec des accents.
% T1 reste pertinent dans un document en anglais contenant ces mots.
%
% Avec LuaLaTeX ou XeLaTeX, utiliser plutôt fontspec :
% ne pas reprendre cette configuration fontenc/inputenc.
\usepackage[T1]{fontenc}

% Règles de césure et intitulés automatiques en anglais.
% Exemple : \tableofcontents porte le titre "Contents".
% Influence également le format de \today.
\usepackage[french]{babel}

% Encodage du FICHIER SOURCE : indique comment lire les caractères
% que tu saisis, par exemple "é".
%
% À distinguer de fontenc :
% - inputenc : lecture des caractères du fichier .tex ;
% - fontenc  : encodage des fontes pour composer le document.
%
% UTF-8 est déjà le défaut des versions récentes de pdfLaTeX :
% cette ligne est donc facultative.
\usepackage[utf8]{inputenc}

% Utilise les fontes vectorielles Latin Modern, une version enrichie
% des fontes Computer Modern traditionnelles de LaTeX.
% Elles restent nettes à tout niveau de zoom et permettent à microtype
% d'ajuster légèrement la largeur des caractères (expansion).
% À charger après \usepackage[T1]{fontenc}.
\usepackage{lmodern}

% Améliore automatiquement la typographie, notamment grâce à de
% légers ajustements de largeur des caractères et à leur placement
% en bord de ligne, selon le moteur et les fontes.
% Aucune commande particulière n'est nécessaire dans le texte.
\usepackage{microtype}

% Désactive la césure automatique des mots.
% Attention : peut entraîner des espaces trop grands entre les mots
% ou des lignes qui débordent lorsque le texte est justifié.
\usepackage[none]{hyphenat}

% Supprime les alinéas et ajoute un espace entre les paragraphes.
\usepackage{parskip}

% Personnalise la mise en forme des listes.
\usepackage{enumitem}

% Aligne les tirets sur la marge gauche.
\setlist[itemize]{
    label=--,
    labelindent=0pt,
    leftmargin=*,
    align=left
}

% --------------------------------------------------
% Mathematics and figures
% --------------------------------------------------

% Outils de composition mathématique : équations alignées,
% matrices, texte dans les formules et références aux équations.
%
% Exemple :
% \begin{align}
%   u &= x + y, \\
%   v &= x - y.
% \end{align}
%
% Exemple de texte dans une formule :
% \[ u = 0 \quad \text{on the boundary}. \]
%
% Exemple de référence : Equation~\eqref{eq:continuity}.
\usepackage{amsmath}

% Symboles mathématiques supplémentaires.
% Charge aussi amsfonts, qui fournit notamment \mathbb.
%
% Exemple : $\mathbf{u} \in \mathbb{R}^{2}$.
% Exemple : $A \subsetneq B$.
\usepackage{amssymb}

% Dessins vectoriels réalisés directement dans le document :
% grilles, axes, flèches et schémas.
%
% Exemple :
% \begin{tikzpicture}
%   \draw[step=1cm, gray] (0,0) grid (3,2);
%   \draw[->] (0,0) -- (3.5,0) node[right] {$x$};
%   \draw[->] (0,0) -- (0,2.5) node[above] {$y$};
% \end{tikzpicture}
\usepackage{tikz}

% Insertion d'images externes avec réglage de taille ou rotation.
% Avec pdfLaTeX, on utilise notamment des PDF, PNG ou JPEG.
%
% Exemple :
% \includegraphics[width=0.7\linewidth]{figures/grid.pdf}
%
% \linewidth désigne la largeur disponible à cet endroit.
\usepackage{graphicx}

% --------------------------------------------------
% Page layout and links
% --------------------------------------------------

% Ajoute « Annexe A », « Annexe B », etc. aux titres et au sommaire.
% Les sections d'annexe sont placées dans l'environnement appendices.
\usepackage[title,titletoc]{appendix}
\renewcommand{\appendixname}{Annexe}

% Définit les marges : ici 2,5 cm sur les quatre côtés.
\usepackage[margin=2.5cm]{geometry}

% Conserve des espacements verticaux naturels, même en recto verso.
% L'espace inutilisé reste en bas de page au lieu d'étirer le contenu.
\raggedbottom

% Numéros en pied de page, côté extérieur pour une reliure à gauche.
% LE : gauche des pages paires ; RO : droite des pages impaires.
\usepackage{fancyhdr}
\pagestyle{fancy}
\fancyhf{}
\fancyfoot[LE,RO]{\thepage}
\renewcommand{\headrulewidth}{0pt}
\renewcommand{\footrulewidth}{0pt}

% Même disposition pour les pages utilisant le style plain.
% Le style empty de la couverture et de la page blanche reste inchangé.
\fancypagestyle{plain}{%
    \fancyhf{}
    \fancyfoot[LE,RO]{\thepage}
    \renewcommand{\headrulewidth}{0pt}
    \renewcommand{\footrulewidth}{0pt}
}

% Gère les références stockées dans bibliography.bib.
% Les citations sont numérotées dans leur ordre d'apparition.
% Exemple dans le texte : \cite{stam1999stable}.
% Latexmk lance Biber pour construire la bibliographie.
\usepackage[backend=biber,style=numeric,sorting=none]{biblatex}
\addbibresource{bibliography.bib}

% Adapte les guillemets à la langue du document.
% Recommandé avec biblatex et babel.
\usepackage{csquotes}

% Rend cliquables les références et les entrées du sommaire.
% Fournit également les commandes \url et \href.
%
% hidelinks masque les couleurs et les cadres des liens,
% mais les liens restent cliquables.
% À charger généralement vers la fin du préambule.
%
% Exemple : \url{https://www.latex-project.org}
% Exemple : \href{https://www.latex-project.org}{LaTeX Project}
\usepackage[hidelinks]{hyperref}

% Met en gras italique les termes cliquables renvoyant à une explication.
% Premier argument : label cible ; second argument : texte affiché.
% Exemple : \conceptlink{annexes:forme_conservative}{forme conservative}
\newcommand{\conceptlink}[2]{%
    \hyperref[#1]{\textbf{\textit{#2}}}%
}



\begin{document}

    \begin{titlepage}
  \centering

  \vspace*{1cm}

  % Chemin à adapter. Privilégie le logo sombre sur une page blanche.
  \includegraphics[width=0.6\textwidth]{../logo/kirin_black_blue.png}

  \vspace{1.5cm}

  {\Huge\bfseries KIRIN\par}

  \vspace{0.5cm}

  {\Large A 2D Incompressible Flow Solver\par}

  \vspace{0.8cm}

  \rule{0.75\textwidth}{0.4pt}

  \vspace{0.8cm}

  {\large From Equations to Code\par}

  \vfill

  {\large Guilhem LAVOINNE\par}

  \vspace{0.4cm}

  {\large \today\par}
\end{titlepage}

    % Page blanche sans en-tête ni numéro après la page de garde.
    \null
    \thispagestyle{empty}
    \newpage

    \tableofcontents
    
    \clearpage

    \section{Introduction}
        KIRIN est un projet pédagogique de développement d'un solveur
bidimensionnel d'écoulements incompressibles, écrit en Rust.
Son objectif principal est de présenter l'ensemble du cheminement
qui conduit des équations physiques à leur implémentation :
choix des hypothèses, formulation du modèle, discrétisation,
construction des algorithmes et traduction en code.

Deux approches numériques seront explorées :
\begin{itemize}
    \item une approche conservative, privilégiant la fidélité
    des résultats et le respect des bilans de conservation ;

    \item une approche inspirée de \textit{Stable Fluids}
    de Jos Stam (1999)~\cite{stam1999stable}, privilégiant la robustesse et l'interactivité.
\end{itemize}

Ce document accompagne cette démarche en explicitant les choix
physiques et numériques, leur mise en œuvre et la vérification
des résultats obtenus.


    \section{Modèle physique}
        \subsection{Les équations de Navier--Stokes}
        Les équations de Navier--Stokes décrivent l'écoulement d'un fluide visqueux.
Les bilans de masse et de quantité de mouvement sous
\conceptlink{annexes:forme_conservative}{forme conservative} sont :
\begin{equation}
    \label{eq:navier_stokes}
    \boxed{
        \begin{aligned}
            \frac{\partial \rho}{\partial t}
            + \nabla \cdot (\rho \mathbf{u})
            &= 0, \\
            \frac{\partial (\rho \mathbf{u})}{\partial t}
            + \nabla \cdot \left(
                \rho \mathbf{u} \otimes \mathbf{u} + p \mathbf{I}
            \right)
            &= \nabla \cdot \boldsymbol{\tau} + \rho \mathbf{g}.
        \end{aligned}}
\end{equation}

La masse volumique est notée $\rho$, la vitesse $\mathbf{u}$,
la pression $p$ et l'accélération de la pesanteur $\mathbf{g}$.
Le tenseur identité est noté $\mathbf{I}$, et $\otimes$ désigne
le produit tensoriel.

Le tenseur $\boldsymbol{\tau}$ représente les contraintes visqueuses.
Pour un \conceptlink{annexes:type_fluids}{fluide newtonien isotrope}, il s'écrit :
\begin{equation}
    \label{eq:viscous_stress}
    \boldsymbol{\tau}
    = \mu \left[\nabla \mathbf{u}
        + (\nabla \mathbf{u})^{\mathsf{T}}\right]
    + \lambda (\nabla \cdot \mathbf{u})\mathbf{I},
\end{equation}
où $\mu$ est la viscosité dynamique de cisaillement et $\lambda$
le second coefficient de viscosité.

Les équations d'Euler correspondent au cas où les contraintes visqueuses
sont négligées : $\boldsymbol{\tau} = \mathbf{0}$.


        
        \subsection{Hypothèses et équations retenues}
        Les équations \eqref{eq:navier_stokes} peuvent être simplifiées sous les
hypothèses suivantes :
\begin{itemize}
    \item l'écoulement est incompressible :
    $\nabla \cdot \mathbf{u} = 0$,

    \item le fluide est homogène et sa masse volumique est constante :
    $\rho(\mathbf{x}, t) = \rho_0 > 0$,

    \item la viscosité dynamique est constante : $\mu = \mu_0$,

    \item la gravité est négligée : $\mathbf{g} = \mathbf{0}$.
\end{itemize}
Pour un fluide newtonien à viscosité constante, la contrainte
d'incompressibilité donne
$\nabla \cdot \boldsymbol{\tau} = \mu_0 \nabla^2 \mathbf{u}$.
En posant $\nu = \mu_0 / \rho_0$, la viscosité cinématique,
on obtient les équations de Navier--Stokes incompressibles :
\begin{equation}
    \label{eq:navier_stokes_incompressible}
    \boxed{
    \begin{aligned}
        \nabla \cdot \mathbf{u} &= 0, \\
        \frac{\partial \mathbf{u}}{\partial t}
        + \nabla \cdot \left(\mathbf{u} \otimes \mathbf{u}\right)
        &= -\frac{1}{\rho_0}\nabla p + \nu \nabla^2 \mathbf{u}.
    \end{aligned}}
\end{equation}

L'écoulement considéré est bidimensionnel :
la vitesse s'écrit $\mathbf{u} = (u,v)$ et les champs
dépendent de $(x,y,t)$, sans dépendance selon $z$.



        \subsection{Transport du colorant}
        Le colorant (\textit{dye} en anglais) est représenté par un champ
scalaire de concentration $c(\mathbf{x}, t)$, transporté par
advection dans le champ de vitesse du fluide.
Il constitue un traceur passif : sa présence n'influence pas
la dynamique de l'écoulement, mais permet d'en visualiser le transport.

En l'absence de diffusion, l'évolution de la concentration satisfait
l'équation suivante sous
\conceptlink{annexes:forme_conservative}{forme conservative} :
\begin{equation}
    \label{eq:dye_transport}
    \boxed{
        \frac{\partial c}{\partial t}
        + \nabla \cdot (c\mathbf{u}) = s.
    }
\end{equation}
Le terme source $s(\mathbf{x}, t)$ représente la quantité de
concentration ajoutée par unité de temps. Ce terme pourra être utilisé pour répondre à
l'injection de colorant par l'utilisateur.

        \subsection{Conditions initiales et aux limites}
        \label{sec:boundary_conditions}

Le domaine de calcul est un rectangle fermé
$\Omega = [0,L_x] \times [0,L_y]$ représenté par les bords de notre simulation, délimité par quatre parois fixes.
Le vecteur $\mathbf{n}$ désigne la normale unitaire sortante à la frontière
$\partial\Omega$.
Les conditions aux limites précisent le comportement de la vitesse,
de la pression et du colorant sur ces parois.

\paragraph*{Vitesse : parois imperméables}

Le fluide ne traverse pas les parois : la composante normale de la vitesse
est nulle sur la frontière.
\begin{equation}
    \label{eq:wall_impermeability}
    \boxed{\mathbf{u} \cdot \mathbf{n} = 0
        \quad \text{sur } \partial\Omega.}
\end{equation}
Avec $\mathbf{u} = (u,v)$, cette condition s'écrit :
\begin{equation}
    \begin{aligned}
        u(0,y,t) = u(L_x,y,t) &= 0, \\
        v(x,0,t) = v(x,L_y,t) &= 0.
    \end{aligned}
\end{equation}

La composante tangentielle n'est pas imposée à zéro : le fluide peut
glisser le long des murs.
Cette condition convient au modèle non visqueux actuellement implémenté.
Les valeurs tangentielles utilisées au bord sont prolongées par symétrie,
sans imposer d'adhérence à la paroi.

Pour une extension visqueuse, le comportement tangentiel devra être
précisé : une paroi fixe avec adhérence impose $\mathbf{u}=\mathbf{0}$,
tandis qu'une paroi à glissement libre impose une contrainte visqueuse
tangentielle nulle.
Sur une paroi plane, avec une viscosité constante non nulle, ce dernier
choix correspond à $\partial u_{\mathrm{t}}/\partial n = 0$,
où $u_{\mathrm{t}}$ est la composante tangentielle de la vitesse.

\paragraph*{Colorant : absence de flux aux parois}

En l'absence de diffusion, le flux de colorant est $c\mathbf{u}$.
L'imperméabilité des murs entraîne donc :
\begin{equation}
    \label{eq:dye_boundary}
    \boxed{(c\mathbf{u})\cdot\mathbf{n} = 0
        \quad \text{sur } \partial\Omega.}
\end{equation}
Aucune concentration nulle n'est imposée au bord : le colorant peut
être présent contre une paroi, mais ne peut pas la traverser.
Dans ce domaine fermé, sa quantité totale ne varie que sous l'effet
du terme source $s$ :
\begin{equation}
    \frac{\mathrm{d}}{\mathrm{d}t}\int_{\Omega} c\,\mathrm{d}\Omega
    = \int_{\Omega} s\,\mathrm{d}\Omega.
\end{equation}

\paragraph*{Conditions initiales}

La configuration initiale de KIRIN correspond à un fluide au repos
et à un domaine sans colorant :
\begin{equation}
    \mathbf{u}(\mathbf{x},0) = \mathbf{0},
    \qquad c(\mathbf{x},0) = 0.
\end{equation}
La pression est initialisée à zéro comme point de départ du solveur
itératif ; elle est ensuite déterminée par la contrainte
d'incompressibilité.
Les interactions de l'utilisateur ajoutent du colorant et modifient
la vitesse à l'intérieur du domaine, sans autoriser de flux à travers
les murs.


    \section{Méthodes de résolution}
        \subsection{Formulation couplée vitesse--pression}
        \subsection{Principe de la méthode de projection}
        \subsection{Principe de l'approche de Stam}
        
    \section{Discrétisation du problème}
        \subsection{Équations en composantes bidimensionnelles}
        \subsection{Maillage et disposition des inconnues}
        \subsection{Opérateurs différentiels discrets}
        \subsection{Discrétisation temporelle}

    \section{Approche conservative avec projection}
        \subsection{Calcul des vitesses intermédiaires}
        \subsection{Résolution de la pression et correction des vitesses}
        \subsection{Transport du colorant}
        \subsection{Contraintes de stabilité}
        \subsection{Implémentation en Rust}
    
    \section{Approche semi-lagrangienne de Stam}
        \subsection{Éléments communs et différences}
        \subsection{Remontée des trajectoires et interpolation}
        \subsection{Stabilité et limites de la méthode}
        \subsection{Implémentation en Rust}

    \section{Conclusion et perspectives}

    \clearpage
    \printbibliography[heading=bibintoc,title={Bibliographie}]
    
    \begin{appendices}
        \clearpage
        \section{Notions et définitions}
\label{annexes:notions}

\subsection{Forme conservative}
\label{annexes:forme_conservative}

Une équation sous forme conservative exprime un bilan local :
la variation d'une quantité par unité de volume résulte des flux
qui la transportent et d'éventuelles sources.
Pour une grandeur scalaire $q$, elle s'écrit :
\begin{equation*}
    \frac{\partial q}{\partial t} + \nabla \cdot \mathbf{F} = s,
\end{equation*}
où $\mathbf{F}$ est le flux et $s$ le terme source.
Sans source ni flux à travers les frontières, la quantité totale
dans un domaine fixe est conservée.
Par exemple, le bilan de masse utilise $q = \rho$,
$\mathbf{F} = \rho\mathbf{u}$ et $s = 0$.

\subsection{Tenseur des contraintes et modèle newtonien isotrope}
\label{annexes:type_fluids}

Le tenseur des contraintes de Cauchy décrit les efforts internes
et se décompose ici en une contribution de pression et une contribution
visqueuse : $\boldsymbol{\sigma} = -p\mathbf{I} + \boldsymbol{\tau}$.
Le modèle newtonien isotrope suppose une relation linéaire entre
les contraintes visqueuses et le tenseur des taux de déformation,
sans direction matérielle privilégiée.
La dérivation de cette relation et les hypothèses associées seront
détaillées ultérieurement.


        \clearpage
        \section{Dérivations et calculs complémentaires}
\label{annexes:details_calculs}

\subsection{Simplification des équations d'Euler}
Cette dérivation concerne le cas non visqueux : l'hypothèse
$\boldsymbol{\tau} = \mathbf{0}$ dans les bilans
\eqref{eq:navier_stokes} donne les équations d'Euler.

La masse volumique est supposée uniforme dans l'espace
et constante dans le temps :
\begin{equation}
    \rho(\mathbf{x}, t) = \rho_0 > 0.
\end{equation}
La gravité est également négligée : $\mathbf{g} = \mathbf{0}$.

\paragraph*{Conservation de la masse}

L'équation de conservation de la masse s'écrit :
\begin{equation}
    \frac{\partial \rho}{\partial t}
    + \nabla \cdot (\rho \mathbf{u}) = 0.
\end{equation}

Le développement de la divergence du produit donne :
\begin{equation}
    \frac{\partial \rho}{\partial t}
    + \mathbf{u} \cdot \nabla \rho
    + \rho \nabla \cdot \mathbf{u} = 0.
\end{equation}

La masse volumique étant constante, ses dérivées temporelle
et spatiales sont nulles :
\begin{equation}
    \frac{\partial \rho}{\partial t} = 0,
    \qquad
    \nabla \rho = \mathbf{0}.
\end{equation}

Il reste donc :
\begin{equation}
    \rho_0 \nabla \cdot \mathbf{u} = 0.
\end{equation}

Puisque $\rho_0 > 0$, il en résulte la contrainte
d'incompressibilité :
\begin{equation}
    \boxed{\nabla \cdot \mathbf{u} = 0.}
\end{equation}

\paragraph*{Conservation de la quantité de mouvement}

L'équation de conservation de la quantité de mouvement s'écrit :
\begin{equation}
    \frac{\partial (\rho \mathbf{u})}{\partial t}
    + \nabla \cdot \left(
        \rho \mathbf{u} \otimes \mathbf{u}
        + p \mathbf{I}
    \right)
    = \rho \mathbf{g}.
\end{equation}

La linéarité de la divergence et l'identité
$\nabla \cdot (p \mathbf{I}) = \nabla p$ donnent :
\begin{equation}
    \frac{\partial (\rho \mathbf{u})}{\partial t}
    + \nabla \cdot
        \left(\rho \mathbf{u} \otimes \mathbf{u}\right)
    + \nabla p
    = \rho \mathbf{g}.
\end{equation}

La masse volumique étant constante et la gravité étant négligée,
cette équation devient :
\begin{equation}
    \rho_0 \frac{\partial \mathbf{u}}{\partial t}
    + \rho_0 \nabla \cdot
        \left(\mathbf{u} \otimes \mathbf{u}\right)
    + \nabla p
    = \mathbf{0}.
\end{equation}

La division par $\rho_0$ et le passage du gradient de pression
dans le membre de droite donnent :
\begin{equation}
    \frac{\partial \mathbf{u}}{\partial t}
    + \nabla \cdot \left(\mathbf{u} \otimes \mathbf{u}\right)
    = -\frac{1}{\rho_0}\nabla p.
\end{equation}

\paragraph*{Système obtenu}

Les deux équations constituent
le système \eqref{eq:euler_simplified} :
\begin{equation*}
    \begin{aligned}
        \nabla \cdot \mathbf{u} &= 0, \\
        \frac{\partial \mathbf{u}}{\partial t}
        + \nabla \cdot \left(\mathbf{u} \otimes \mathbf{u}\right)
        &= -\frac{1}{\rho_0}\nabla p.
    \end{aligned}
\end{equation*}


    \end{appendices}


\end{document}
